% supplementary.tex -- standalone compilable Supplementary Material for
% BSPC. Not \input from paper.tex; compiled to its own PDF, as is standard
% for Elsevier supplementary submissions.
%
% Scope, per direct instruction: full figures, full tables, implementation
% detail, robustness, and extra statistics only. No result already in the
% main text is re-explained or re-interpreted here -- captions state what is
% shown, not what it means; that reading is in the main manuscript.

\documentclass{article}

\usepackage[margin=2.5cm]{geometry}
\usepackage{amsmath}
\usepackage{amssymb}
\usepackage{graphicx}
\usepackage{booktabs}
\usepackage{multirow}
\usepackage{xr}
\usepackage{hyperref}

% Cross-references into the main manuscript's section/figure/table numbers.
% Requires paper.tex to have been compiled at least once (paper.aux present)
% -- if numbers show as "??", compile paper.tex first, then this file.
\externaldocument{paper}

\renewcommand{\thefigure}{S\arabic{figure}}
\renewcommand{\thetable}{S\arabic{table}}
\renewcommand{\thesection}{S}

% \suppsec{<n>}{<title>} -- open supplementary section S<n> AND pin the figure
% and table counters to it, so every float inside the section is numbered S<n>.
% Heading number and float numbers are driven by the same argument and cannot
% diverge. They previously did: with plain \subsection* headings and free-running
% counters, section "S4" contained a float LaTeX numbered "S3" and section "S7"
% contained one numbered "S4", while the main text referred to the latter, by
% hardcoded string, as "Supplementary Table S7". The main text deliberately does
% not \usepackage{xr} into this file (an arXiv build of paper.tex alone would
% then render every such reference as "??"), so those strings cannot be
% link-checked by LaTeX -- which is exactly why the numbers they name have to be
% guaranteed here instead.
%
% Invariant: at most one figure and at most one table per section. A section
% needing two tables is split into two sections (see S8/S9), rather than being
% allowed to introduce a second, unnumbered float.
\newcommand{\suppsec}[2]{%
  \setcounter{figure}{#1}\addtocounter{figure}{-1}%
  \setcounter{table}{#1}\addtocounter{table}{-1}%
  \subsection*{S#1. #2}%
}

\title{Supplementary Material\\[0.3em]
\large Decodable but Not Accessible: Auditing Distance-Based Reliability Estimation on Disentangled Skin-Lesion Representations}
\author{}
\date{}

\begin{document}
\maketitle
\vspace{-1.5em}

\section*{Supplementary Material}

\suppsec{1}{Full geometry metrics across the ladder}

Extends main-text Figure~\ref{fig:ladder-trend}, which shows only condition number and the decoupled Fisher-ratio scalar. Figure~\ref{fig:s1} shows all five pre-specified geometry metrics (condition number, Fisher ratio---Hotelling--Lawley and decoupled-scalar forms---and Mardia's multivariate kurtosis, $b$ and $z$ statistics), one gray point per checkpoint ($n=13$), mean $\pm$ SD in blue, across the three rungs of the disentanglement ladder. Table~\ref{tab:s1-geometry-stats} reports the corresponding exact-permutation Kendall's $\tau$ and Jonckheere--Terpstra statistics for all five metrics.

\begin{figure}[htbp]
\centering
\includegraphics[width=\textwidth]{figures/figure1_e1_ladder_trend.png}
\caption{All five geometry metrics across the disentanglement ladder.}
\label{fig:s1}
\end{figure}

% table_s1_geometry_stats.tex -- from results/e1_kendall_tau.csv and
% results/e1_jonckheere_terpstra.csv. Full 5-metric statistics for geometry
% vs. rung order (main text Table~\ref{tab:geometry-summary} reports only
% condition number's headline statistic).
\begin{table}[htbp]
\centering
\caption{Geometry vs.\ rung order: full statistics for all five metrics (exact-permutation tests, Section~\ref{sec:stats}).}
\label{tab:s1-geometry-stats}
\small
\begin{tabular}{lcccccc}
\toprule
Metric & $\tau$ (full, $n=13$) & $p$ (exact) & $\tau$ (common-seed, $n=9$) & $p$ (exact) & $J$ & $p_{JT}$ (exact) \\
\midrule
Condition number             & $0.840$  & $2.8\times10^{-5}$ & $0.866$  & $0.0012$ & $55.0$ & $1.4\times10^{-5}$ \\
Fisher ratio (HL)            & $0.168$  & $0.521$            & $-0.096$ & $0.831$  & $33.0$ & $0.261$ \\
Fisher ratio (scalar)        & $0.076$  & $0.798$            & $-0.289$ & $0.388$  & $30.0$ & $0.399$ \\
Mardia's kurtosis ($b$)      & $-0.137$ & $0.608$            & $-0.160$ & $0.668$  & $23.0$ & $0.739$ \\
Mardia's kurtosis ($z$)      & $-0.168$ & $0.521$            & $-0.160$ & $0.668$  & $22.0$ & $0.780$ \\
\bottomrule
\end{tabular}
\end{table}


\suppsec{2}{Full geometry-vs-AUROC statistics}

Extends main-text Figure~\ref{fig:geometry-vs-auroc}, which shows only the condition-number panel. Figure~\ref{fig:s2} shows all five geometry metrics plotted against Mahalanobis AUROC, one point per checkpoint, colored by rung. Table~\ref{tab:s2-auroc-stats} reports the corresponding statistics, including a $10^5$-draw Monte Carlo cross-check of the exact Kendall's $\tau$ $p$-values.

\begin{figure}[htbp]
\centering
\includegraphics[width=\textwidth]{figures/figure2_e2_geometry_vs_auroc.png}
\caption{All five geometry metrics vs.\ Mahalanobis AUROC.}
\label{fig:s2}
\end{figure}

% table_s2_auroc_stats.tex -- from results/e2_kendall_tau.csv and
% results/e2_jonckheere_terpstra_auroc.csv. Full 5-metric statistics for
% geometry vs. Mahalanobis AUROC (main text Figure~\ref{fig:geometry-vs-auroc}
% shows only the condition-number panel).
\begin{table}[htbp]
\centering
\caption{Geometry vs.\ Mahalanobis AUROC: full statistics for all five metrics ($n=13$, exact-permutation Kendall's $\tau$; Monte Carlo cross-check at $10^5$ draws in parentheses).}
\label{tab:s2-auroc-stats}
\small
\begin{tabular}{lcc}
\toprule
Metric & $\tau$ & $p$ (exact / MC) \\
\midrule
Condition number       & $-0.128$ & $0.590$ / $0.589$ \\
Fisher ratio (HL)      & $-0.282$ & $0.204$ / $0.202$ \\
Fisher ratio (scalar)  & $-0.154$ & $0.510$ / $0.509$ \\
Mardia's kurtosis ($b$) & $-0.128$ & $0.590$ / $0.590$ \\
Mardia's kurtosis ($z$) & $-0.128$ & $0.590$ / $0.592$ \\
\bottomrule
\end{tabular}

\vspace{0.3cm}
\small AUROC itself vs.\ rung order (Jonckheere--Terpstra): $J=25.0$ (null mean $27.5$), $p_{\mathrm{exact}}=0.650$.
\end{table}


\suppsec{3}{Baseline reference: implementation detail}

The \texttt{baseline\_soft} checkpoint (main text, Section~\ref{sec:results-baseline}) is reported as a descriptive reference, not a fourth ladder point. Table~\ref{tab:s3-baseline} gives its full distance-scorer and probe AUROC (single checkpoint, seed 42). Its geometry metrics (condition number, Fisher ratio, Mardia's kurtosis) were not computed: because \texttt{baseline\_soft} was never load-bearing for any statistical claim in this study, computing its geometry was out of scope.

\begin{table}[htbp]
\centering
\caption{\texttt{baseline\_soft}: distance-scorer and domain-probe AUROC ($n=1$ checkpoint).}
\label{tab:s3-baseline}
\begin{tabular}{lc}
\toprule
Scorer / probe & AUROC \\
\midrule
Mahalanobis          & 0.863 \\
Cosine-to-centroid   & 0.715 \\
$k$-NN ($k=1$)       & 0.865 \\
$k$-NN ($k=10$)      & 0.853 \\
$k$-NN ($k=50$)      & 0.844 \\
Distance-scorer mean & 0.828 \\
\midrule
Probe: logistic regression & 0.998 \\
Probe: linear SVM          & 0.997 \\
Probe: random forest       & 0.977 \\
\bottomrule
\end{tabular}
\end{table}

\suppsec{4}{ID/OOD Mahalanobis distance distributions}

Figure~\ref{fig:s4} shows the full histogram, kernel density estimate, and empirical CDF of the ID (ISIC-test) and OOD (PAD-UFES) Mahalanobis squared-distance distributions, pooled across seeds within each rung, supporting the median-gap numbers reported in main-text Section~\ref{sec:results-distance-gap}. The distributions are visibly overlapping, not cleanly separated.

\begin{figure}[htbp]
\centering
\includegraphics[width=\textwidth]{figures/figure_e2_5_distance_distributions.png}
\caption{ID vs.\ OOD Mahalanobis squared-distance distributions, all three rungs.}
\label{fig:s4}
\end{figure}

\suppsec{5}{Full nearest-centroid confusion matrix}

Figure~\ref{fig:s5} shows the complete $8\times8$ nearest-centroid confusion matrix for ISIC-test samples (true label vs.\ predicted centroid), supporting the majority-class-only summary in main-text Figure~\ref{fig:mechanism}B.

\begin{figure}[htbp]
\centering
\includegraphics[width=0.85\textwidth]{figures/figure_id_confusion_matrix.png}
\caption{Full ISIC-test nearest-centroid confusion matrix, all 8 classes.}
\label{fig:s5}
\end{figure}

\suppsec{6}{Full predicted-class distribution for PAD-UFES}

Figure~\ref{fig:s6} shows the complete nearest-centroid predicted-class distribution for PAD-UFES samples across all 8 ISIC classes, extending main-text Figure~\ref{fig:mechanism}B's majority-class (Nevus)-only summary.

\begin{figure}[htbp]
\centering
\includegraphics[width=0.85\textwidth]{figures/figure_predicted_class_ood.png}
\caption{Full PAD-UFES nearest-centroid predicted-class distribution, all 8 classes.}
\label{fig:s6}
\end{figure}

\suppsec{7}{Full statistical test results: scorers and probes vs.\ rung order}

Table~\ref{tab:s7-scorer-probe-stats} reports exact-permutation Kendall's $\tau$ (and, for the five distance-based scorers, Jonckheere--Terpstra) for AUROC vs.\ rung order, for every scorer variant (main-text Section~\ref{sec:results-scorer-comparison}) and every probe (main-text Section~\ref{sec:results-probe}) individually. No metric or scorer reaches significance at $n=13$; none was omitted from this table.

% table_s7_scorer_probe_stats.tex -- from results/e2_6_kendall_tau.csv,
% results/e2_6_jonckheere_terpstra.csv, results/e2_7_kendall_tau.csv.
% Full trend statistics for all eight scorers (five distance-based, three
% non-distance-based added in E2.8) and all three domain probes, vs. rung
% order (main text reports only the range of p-values in prose, plus
% Energy's isolated trend discussed explicitly, Sections~
% \ref{sec:results-scorer-comparison}--\ref{sec:results-probe}).
\begin{table}[htbp]
\centering
\caption{AUROC vs.\ rung order, per scorer and per probe ($n=13$, exact-permutation tests).}
\label{tab:s7-scorer-probe-stats}
\small
\begin{tabular}{llccc}
\toprule
& Scorer / probe & $\tau$ & $p$ (Kendall, exact) & $p_{JT}$ (exact) \\
\midrule
\multirow{5}{*}{Distance-based}
& Mahalanobis        & $-0.076$ & $0.798$ & $0.650$ \\
& Cosine-to-centroid & $0.321$  & $0.194$ & $0.097$ \\
& $k$-NN ($k=1$)     & $-0.046$ & $0.898$ & $0.601$ \\
& $k$-NN ($k=10$)    & $0.015$  & $1.000$ & $0.500$ \\
& $k$-NN ($k=50$)    & $0.076$  & $0.798$ & $0.399$ \\
\midrule
\multirow{3}{*}{Non-distance-based}
& Energy             & $0.443$  & $0.065$ & $0.032$ \\
& ViM                & $0.076$  & $0.798$ & $0.399$ \\
& Density (KDE)      & $0.015$  & $1.000$ & $0.500$ \\
\midrule
\multirow{3}{*}{Domain probes}
& Logistic regression & $0.260$  & $0.302$ & --- \\
& Linear SVM          & $0.260$  & $0.302$ & --- \\
& Random forest       & $-0.107$ & $0.701$ & --- \\
\bottomrule
\end{tabular}
\end{table}


\suppsec{8}{Detectability and power of the permutation designs}

Supports main-text Section~\ref{sec:power} and the Limitations discussion in Section~\ref{sec:discussion-limitations}. Figure~\ref{fig:s8}A shows the power of each permutation design as a function of the expected Kendall's $\tau$ under the alternative, with the design's critical value $\tau_{\mathrm{crit}}$ and the pre-specified $|\tau|\geq0.3$ threshold marked; Figure~\ref{fig:s8}B shows every observed $\tau$ with its 95\% bootstrap interval against the band of values no design at this sample size can resolve. Table~\ref{tab:s8-power} gives the design-level numbers; the per-test intervals plotted in panel B are tabulated in Section~S9.

\begin{figure}[htbp]
\centering
\includegraphics[width=\textwidth]{figures/figure_s_power.png}
\caption{(A) Power of each exact permutation design against expected $\tau$ ($2\times10^{4}$ simulated datasets per point, common random numbers; $\alpha=0.05$, two-sided). Dotted vertical lines mark each design's $\tau_{\mathrm{crit}}$; the dashed red line marks the pre-specified $|\tau|\geq0.3$ criterion, which falls below every $\tau_{\mathrm{crit}}$. (B) Observed $\tau$ and 95\% BCa bootstrap intervals for all 21 reported associations. The shaded band spans $\pm\tau_{\mathrm{crit}}$ for the 5/5/3 ladder: no $|\tau|$ inside it can reach $p\leq0.05$ at this design. Condition number vs.\ $\lambda_{orth}$ sits at the design ceiling and is shown as a point without an interval (see Table~\ref{tab:s9-tau-ci}).}
\label{fig:s8}
\end{figure}

% table_s8_power.tex -- from results/power_design_summary.csv
% (analysis/analyze_power_and_ci.py). Design-level detectability: what each
% of this study's two permutation designs could have found, independent of
% the data actually observed.
\begin{table}[htbp]
\centering
\caption{Detectability of the two permutation designs used in this study. $\tau_{\mathrm{crit}}$ is the smallest $|\tau|$ attainable at that design whose two-sided exact permutation $p$-value reaches $0.05$: no dataset, however clean, can produce a significant association below it. Power is computed by simulation ($2\times10^{4}$ datasets per point, common random numbers) under an alternative whose expected value of the reported statistic equals the stated $\tau$ --- a bivariate-normal alternative for Design~A, a linear dose-shift alternative $y=\theta d+\varepsilon$ for Design~B. MDE$_{80\%}$ is the expected $\tau$ at which power first reaches $0.80$.}
\label{tab:s8-power}
\small
\begin{tabular}{lccccccc}
\toprule
& & Permutation & Max.\ attainable & & \multicolumn{2}{c}{Power} & \\
\cmidrule(lr){6-7}
Design & $n$ & null space & $|\tau|$ & $\tau_{\mathrm{crit}}$ & $\tau=0.3$ & $\tau=0.5$ & MDE$_{80\%}$ \\
\midrule
A: continuous--continuous            & 13 & $13!$ orderings & 1.000 & 0.436 & 0.27 & 0.71 & 0.54 \\
\quad (E2: geometry vs.\ AUROC)      &    &                 &       &       &      &      &      \\
B: 5/5/3 ordered ladder              & 13 & 72,072          & 0.840 & 0.473 & 0.23 & 0.63 & 0.57 \\
\quad (E1, E2.6, E2.7 vs.\ $\lambda_{orth}$) &  &            &       &       &      &      &      \\
B: 3/3/3 common-seed subset          &  9 & 1,680           & 0.866 & 0.609 & 0.14 & 0.39 & 0.69 \\
\bottomrule
\end{tabular}

\vspace{0.25cm}
\small The success criterion this study pre-specified (Section~\ref{sec:stats}, $|\tau|\geq0.3$ \emph{and} significant at $\alpha=0.05$) lies below $\tau_{\mathrm{crit}}$ for all three designs, and is therefore unsatisfiable at these sample sizes: an association of exactly the pre-specified magnitude could not have been declared significant regardless of the data. Only the condition-number effect ($\tau=0.840$, the maximum attainable value at this design) clears $\tau_{\mathrm{crit}}$.
\end{table}


\suppsec{9}{Interval estimates for every reported association}

Table~\ref{tab:s9-tau-ci} tabulates the point estimates and intervals plotted in Figure~\ref{fig:s8}B: every Kendall's $\tau$ reported anywhere in this study, with a 95\% stratified-bootstrap BCa interval and the largest effect size that interval leaves unexcluded. No association tested is omitted.

% table_s9_tau_ci.tex -- from results/kendall_tau_ci.csv
% (analysis/analyze_power_and_ci.py). Every Kendall's tau reported in this
% paper, with a 95% stratified-bootstrap BCa interval, so that each null
% result states which effect sizes it does and does not exclude rather than
% only that p > 0.05.
\begin{table}[htbp]
\centering
\caption{Every reported Kendall's $\tau$ with a 95\% bootstrap confidence interval ($2\times10^{4}$ stratified resamples, BCa; resampling is within-rung so the 5/5/3 design is held fixed). ``Largest $|\tau|$ not excluded'' is the interval endpoint furthest from zero --- the effect size each null result still permits.}
\label{tab:s9-tau-ci}
\small
\begin{tabular}{llccc}
\toprule
& Test & $\tau$ & 95\% CI & Largest $|\tau|$ not excluded \\
\midrule
\multirow{5}{*}{E1 vs.\ $\lambda_{orth}$}
& Condition number        & $0.840$  & \multicolumn{2}{c}{at design ceiling --- see note} \\
& Fisher ratio (HL)       & $0.168$  & $[-0.461,\ 0.620]$ & 0.62 \\
& Fisher ratio (scalar)   & $0.076$  & $[-0.532,\ 0.532]$ & 0.53 \\
& Mardia's kurtosis ($b$) & $-0.137$ & $[-0.400,\ 0.206]$ & 0.40 \\
& Mardia's kurtosis ($z$) & $-0.168$ & $[-0.412,\ 0.174]$ & 0.41 \\
\midrule
\multirow{5}{*}{E2 vs.\ AUROC}
& Condition number        & $-0.128$ & $[-0.507,\ 0.200]$ & 0.51 \\
& Fisher ratio (HL)       & $-0.282$ & $[-0.750,\ 0.183]$ & 0.75 \\
& Fisher ratio (scalar)   & $-0.154$ & $[-0.611,\ 0.324]$ & 0.61 \\
& Mardia's kurtosis ($b$) & $-0.128$ & $[-0.500,\ 0.233]$ & 0.50 \\
& Mardia's kurtosis ($z$) & $-0.128$ & $[-0.534,\ 0.216]$ & 0.53 \\
\midrule
\multirow{8}{*}{E2.6 scorers}
& Mahalanobis             & $-0.076$ & $[-0.524,\ 0.432]$ & 0.52 \\
& Cosine-to-centroid      & $0.321$  & $[-0.210,\ 0.688]$ & 0.69 \\
& $k$-NN ($k=1$)          & $-0.046$ & $[-0.500,\ 0.556]$ & 0.56 \\
& $k$-NN ($k=10$)         & $0.015$  & $[-0.564,\ 0.489]$ & 0.56 \\
& $k$-NN ($k=50$)         & $0.076$  & $[-0.524,\ 0.545]$ & 0.55 \\
& Energy                  & $0.443$  & $[\ \ 0.048,\ 0.747]$ & 0.75 \\
& ViM                     & $0.076$  & $[-0.268,\ 0.270]$ & \textbf{0.27} \\
& Density (KDE)           & $0.015$  & $[-0.560,\ 0.471]$ & 0.56 \\
\midrule
\multirow{3}{*}{E2.7 probes}
& Logistic regression     & $0.260$  & $[-0.210,\ 0.656]$ & 0.66 \\
& Linear SVM              & $0.260$  & $[-0.210,\ 0.656]$ & 0.66 \\
& Random forest           & $-0.107$ & $[-0.560,\ 0.368]$ & 0.56 \\
\bottomrule
\end{tabular}

\vspace{0.25cm}
\small \emph{Note.} Condition number vs.\ $\lambda_{orth}$ attains $\tau=0.840$, the maximum value the 5/5/3 design permits (a perfectly monotone ladder). Every bootstrap resample is therefore at or below the observed value and the resampling distribution is one-sided by construction, so no bootstrap interval is valid there; the exact permutation $p$-value ($2.8\times10^{-5}$) is the meaningful inference for that effect. ViM (bold) is the only null result in this study whose interval excludes the pre-specified $|\tau|=0.3$; the other nineteen do not.
\end{table}


\end{document}
